\documentclass[12pt,twoside]{report}

% some definitions for the title page
\newcommand{\reporttitle}{The Emotion Engine}
\newcommand{\reportauthor}{Kaede Sugano}
\newcommand{\supervisor}{William Knottenbelt}
\newcommand{\reporttype}{Undergraduate Individual Project}
\newcommand{\degreetype}{MEng} 

% load some definitions and default packages
\input{includes}

% load some macros
\input{notation}

% load title page
\begin{document}
\input{titlepage}


% page numbering etc.
\pagenumbering{roman}
\clearpage{\pagestyle{empty}\cleardoublepage}
\setcounter{page}{1}
\pagestyle{fancy}
 
%%%%%%%%%%%%%%%%%%%%%%%%%%%%%%%%%%%%
\chapter{Implementation and Preliminary Experiments}

This chapter describes the first phase of implementation, which focuses on adapting the Coconut latent reasoning architecture to the affective regression task and diagnosing the representational capacity of the resulting latent space.

\section{Dataset}

All experiments use the \textbf{EmoBank} corpus~\cite{buechel-hahn-2017-emobank}, a sentence-level dataset annotated with continuous Pleasure--Arousal--Dominance (PAD) scores on a scale from 1 to 5.
The dataset was split into 7{,}801 training samples and 861 test samples.
Each training instance consists of an input sentence paired with a gold PAD triple $\{V, A, D\} \in [1,5]^3$.
Reasoning-step annotations were generated offline using GPT-4o and cached as a \texttt{JSONL} file, providing a chain-of-thought scaffold for the staged training curriculum described below.

\section{Model Variants}

Four model variants were trained and evaluated, all using GPT-2 (117M parameters)~\cite{radford2019language} as the backbone.

\paragraph{CoT baseline.}
A standard chain-of-thought model in which each training example is formatted as a free-text reasoning trace followed by the gold PAD answer.
The model is supervised with cross-entropy on the full token sequence.

\paragraph{Coconut baseline.}
The original Coconut architecture~\cite{hao2025traininglargelanguagemodels} replaces textual reasoning tokens with a fixed number of \emph{latent tokens}.
At each training stage $k$, the first $k$ reasoning steps are replaced by latent embeddings that are filled in with the model's own hidden states in a multi-pass forward procedure.
The number of latent tokens per stage is controlled by the hyperparameter $c_{\text{thought}} = 2$, and up to $k_{\max} = 5$ stages are scheduled.

\paragraph{CoconutGPT-SimCoT.}
Extends the Coconut baseline with an auxiliary explainability head that, during training only, decodes each latent reasoning step into a natural-language phrase.
This is analogous to the SIM-CoT supervision strategy~\cite{wei2025simcot} and is intended to prevent latent collapse by ensuring each latent vector encodes semantically distinct content.

\paragraph{CoconutGPT-Factored.}
A new architecture introduced in this project.
Two auxiliary losses are added on top of the Coconut latent pass:
\begin{enumerate}
    \item \textbf{VAD regression loss.} A linear head $\mathbf{W}_{\text{vad}} \in \mathbb{R}^{3 \times d}$ maps the mean-pooled latent representations to a predicted PAD triple $\hat{\mathbf{v}} \in \mathbb{R}^3$.
    The regression loss is:
    \[
        \mathcal{L}_{\text{reg}} = \frac{1}{3} \|\hat{\mathbf{v}} - \mathbf{v}^*\|_2^2
    \]
    where $\mathbf{v}^* = (V^*, A^*, D^*)$ is the gold annotation.

    \item \textbf{Contrastive geometric loss.} Within each mini-batch of size $B$, the cosine similarity matrix of mean-pooled latent vectors is aligned to a target similarity derived from PAD-space Euclidean distance:
    \[
        s_{ij}^* = \exp\!\left(-\frac{\|\mathbf{v}_i^* - \mathbf{v}_j^*\|_2}{\tau}\right), \quad i \neq j
    \]
    The contrastive loss is the mean squared error between the computed cosine similarities and $s_{ij}^*$ over all off-diagonal pairs.
\end{enumerate}

The total training objective is:
\[
    \mathcal{L} = \mathcal{L}_{\text{LM}} + \lambda_{\text{reg}}\,\mathcal{L}_{\text{reg}} + \lambda_{\text{con}}\,\mathcal{L}_{\text{con}}
\]
with $\lambda_{\text{reg}} = 1.0$, $\lambda_{\text{con}} = 0.1$, and $\tau = 1.0$.
The model is warm-started from the Coconut baseline checkpoint.

\section{Phase 1: Linear Probe Diagnosis}

Before training CoconutGPT-Factored, we conducted a linear probing experiment to determine whether PAD information was already linearly decodable from the Coconut baseline's latent representations.

For each sample, a forward pass was performed through the Coconut baseline (checkpoint at epoch 25).
The mean-pooled hidden states at all latent token positions were extracted as a 768-dimensional vector.
A Ridge regression probe was then trained on the 7{,}801 training representations and evaluated on the 861 test representations independently for each PAD dimension.

\begin{table}[H]
\centering
\begin{tabular}{lccc}
\toprule
\textbf{Dimension} & \textbf{Pearson $r$ (test)} & \textbf{$R^2$ (test)} & \textbf{$R^2$ (train)} \\
\midrule
Valence (V)   & 0.695 & \phantom{$-$}0.474 & 0.860 \\
Arousal (A)   & 0.371 & $-$0.000 & 0.804 \\
Dominance (D) & 0.409 & \phantom{$-$}0.101 & 0.676 \\
\bottomrule
\end{tabular}
\caption{Linear probe results on the Coconut baseline latent representations. Arousal shows near-zero test $R^2$, indicating that the information is not linearly encoded in the baseline latent space.}
\label{tab:probe_baseline}
\end{table}

The results reveal a clear asymmetry: Valence is moderately recoverable (Pearson $r = 0.695$), consistent with the hypothesis that surface-level lexical sentiment is already encoded in the latent representations.
By contrast, Arousal achieves a test $R^2$ of $-0.000$, which is below the constant-mean baseline, and Dominance achieves only $R^2 = 0.101$.
The large gap between training and test $R^2$ across all dimensions indicates that the latent space is high-dimensional relative to the amount of affective information it encodes, leading to over-fitting in the probe even with regularisation.

This diagnostic confirms the central hypothesis motivating CoconutGPT-Factored: \textbf{the Coconut baseline does not encode Arousal or Dominance in its latent representations}, which directly explains its inferior performance on those dimensions relative to the CoT model.

\section{Evaluation Metrics}

All models are evaluated on the 861-sample test split using the following metrics:

\begin{itemize}
    \item \textbf{Pearson $r$} per PAD dimension, measuring linear correlation between predicted and gold values.
    \item \textbf{Tolerance accuracy}: the proportion of test samples for which all three predicted dimensions fall within $\pm 0.25$ of the gold values.
    \item \textbf{MAE} and \textbf{RMSE} over all three dimensions.
\end{itemize}

Predicted PAD values are extracted by parsing the model's generated output string.

\section{Results}

Table~\ref{tab:main_results} summarises the evaluation results for all four models.

\begin{table}[H]
\centering
\begin{tabular}{lcccc}
\toprule
\textbf{Model} & \textbf{Pearson V} & \textbf{Pearson A} & \textbf{Pearson D} & \textbf{Tol-Acc} \\
\midrule
CoT baseline             & 0.6858 & 0.4494 & 0.4548 & 46.92\% \\
Coconut baseline         & 0.6841 & 0.3417 & 0.3363 & 47.74\% \\
CoconutGPT-SimCoT        & 0.6768 & 0.2910 & 0.3261 & 42.86\% \\
\textbf{CoconutGPT-Factored} & \textbf{0.8466} & \textbf{0.8119} & \textbf{0.8411} & 45.30\% \\
\bottomrule
\end{tabular}
\caption{Pearson correlation and tolerance accuracy on the EmoBank test set. CoconutGPT-Factored substantially outperforms all baselines on Arousal and Dominance.}
\label{tab:main_results}
\end{table}

CoconutGPT-Factored achieves a substantial improvement on Arousal (0.34 $\to$ 0.81) and Dominance (0.34 $\to$ 0.84), confirming that the VAD regression loss successfully drives the model to encode these dimensions.
Valence also improves from 0.68 to 0.85.
Tolerance accuracy (0.453) is slightly lower than the Coconut baseline (0.477), reflecting the fact that absolute numeric precision is a stricter criterion than linear correlation.

A follow-up linear probe on the CoconutGPT-Factored representations (Table~\ref{tab:probe_factored}) reveals that, while Valence probe accuracy improves slightly ($r = 0.695 \to 0.732$), Arousal and Dominance probes remain weak ($r \approx 0.37$).
This apparent contradiction -- strong end-to-end performance but weak linear probe -- indicates that the VAD information is encoded in a form that is accessible to the \texttt{vad\_head} but not to a purely linear probe trained on question-only contexts.
The discrepancy arises because the probe operates on latent vectors extracted before the answer prefix, whereas the \texttt{vad\_head} at inference time has access to the full context including the answer preamble.
This is an important limitation that motivates future work on aligning the latent space geometry with the PAD axes more directly.

\begin{table}[H]
\centering
\begin{tabular}{lccc}
\toprule
\textbf{Dimension} & \textbf{Pearson $r$ (test)} & \textbf{$R^2$ (test)} & \textbf{$R^2$ (train)} \\
\midrule
Valence (V)   & 0.732 & 0.535 & 0.972 \\
Arousal (A)   & 0.366 & 0.025 & 0.946 \\
Dominance (D) & 0.383 & 0.086 & 0.931 \\
\bottomrule
\end{tabular}
\caption{Linear probe results on CoconutGPT-Factored latent representations. Despite strong end-to-end performance, Arousal and Dominance remain weakly decodable by a linear probe on question-only contexts.}
\label{tab:probe_factored}
\end{table}

\section{Discussion}

The results of this first experimental phase establish two key findings.

First, the linear probe confirms that the Coconut baseline's latent space does not encode Arousal or Dominance, which explains why it underperforms the CoT model despite having access to additional computation via continuous latent tokens.
The latent tokens in the Coconut baseline appear to serve primarily as a computation budget for Valence-related reasoning, which is strongly correlated with surface-level lexical sentiment.

Second, the addition of direct VAD supervision via the regression and contrastive losses in CoconutGPT-Factored yields dramatic improvements on Arousal and Dominance, surpassing the CoT baseline by a large margin.
This demonstrates that continuous latent space reasoning, when appropriately supervised, can encode richer affective structure than token-level generation alone.

Nevertheless, the weak linear probe results on the Factored model suggest that the improvement is driven by the \texttt{vad\_head} decoding during generation, rather than by a fundamental reorganisation of the latent geometry.
Achieving a latent space that is intrinsically aligned with PAD axes -- such that the latent vectors themselves constitute interpretable, steerable affective representations -- remains an open challenge and constitutes the primary direction for the next phase of this project.

%%%%%%%%%%%%%%%%%%%%%%%%%%%%%%%%%%%%
\chapter{Phase 2: Towards Steerable Latent Geometry}

The linear probe results from Phase~1 establish a clear gap: although CoconutGPT-Factored achieves high end-to-end Pearson correlation (A: 0.81, D: 0.84), the Arousal and Dominance dimensions remain weakly decodable from the latent representations by a linear probe ($r \approx 0.37$).
This matters because steering vectors are inherently linear interventions.
If A and D are not linearly encoded in the latent space, a steering direction $\mathbf{d}_A$ derived from contrastive pairs cannot reliably shift Arousal without also disturbing Valence or semantic content.
This chapter describes the architectural and training modifications introduced in Phase~2 to address this limitation.

\section{Diagnosis: Why Does the Probe Fail?}

The Phase~1 architecture uses a single shared regression head $\mathbf{W}_{\text{vad}} \in \mathbb{R}^{3 \times d}$, whose three rows correspond to V, A, and D respectively.
During optimisation, gradients from all three dimensions flow through the same weight matrix.
Because Valence correlates strongly with surface lexical sentiment -- a signal the model acquires readily -- the shared head is dominated by the V gradient, leaving A and D directions weakly constrained.
The contrastive loss similarly uses combined PAD-space Euclidean distance $\|\mathbf{v}_i^* - \mathbf{v}_j^*\|_2$, which conflates all three dimensions and further biases the latent geometry toward Valence.

A second contributing factor is the large gap between $R^2_{\text{train}}$ and $R^2_{\text{test}}$ (e.g.\ A: 0.946 vs.\ 0.025).
This indicates that the linear probe overfits to training patterns rather than recovering a stable geometric direction in the latent space, consistent with the latent representations not having been explicitly regularised to place A and D along separable linear axes.

\section{Architectural Changes}

\subsection{Separate Per-Dimension Projection Heads}

The single shared head $\mathbf{W}_{\text{vad}}$ is replaced by three independent scalar projection heads:
\[
    \hat{V} = \mathbf{w}_V^\top \mathbf{z}, \quad
    \hat{A} = \mathbf{w}_A^\top \mathbf{z}, \quad
    \hat{D} = \mathbf{w}_D^\top \mathbf{z}
\]
where $\mathbf{w}_V, \mathbf{w}_A, \mathbf{w}_D \in \mathbb{R}^d$ are fully independent parameter vectors and $\mathbf{z}$ is the mean-pooled latent representation.
Each head receives its own gradient signal, preventing the V gradient from dominating the update to A and D directions.

\subsection{Dimension-Weighted Regression Loss}

The regression loss is reformulated as a weighted sum over dimensions:
\[
    \mathcal{L}_{\text{reg}} = \sum_{k \in \{V, A, D\}} w_k \cdot \frac{1}{B}\sum_{i=1}^{B}(\hat{k}_i - k_i^*)^2
\]
with weights $(w_V, w_A, w_D) = (1.0, 2.0, 2.0)$, upweighting A and D to compensate for the natural advantage of Valence.

\subsection{Per-Dimension Contrastive Loss}

Rather than computing a single contrastive target from the combined PAD distance, the contrastive loss is decomposed into three independent terms, one per dimension:
\[
    s_{ij}^{*(k)} = \exp\!\left(-\frac{|k_i^* - k_j^*|}{\tau}\right), \quad k \in \{V, A, D\}
\]
\[
    \mathcal{L}_{\text{con}} = \sum_{k \in \{V, A, D\}} w_k \cdot \frac{1}{|\mathcal{P}|}\sum_{(i,j) \in \mathcal{P}} \bigl(\cos(\mathbf{z}_i, \mathbf{z}_j) - s_{ij}^{*(k)}\bigr)^2
\]
where $\mathcal{P}$ denotes all off-diagonal pairs in the mini-batch.
This formulation pulls the latent space geometry toward each PAD axis independently, weighted by the same $(1.0, 2.0, 2.0)$ coefficients.

\subsection{Orthogonality Regularisation}

To encourage V, A, and D to occupy geometrically distinct directions in the latent space, an orthogonality penalty is imposed on the three head weight vectors:
\[
    \mathcal{L}_{\text{orth}} = \left\| \hat{\mathbf{W}}\hat{\mathbf{W}}^\top - \mathbf{I}_3 \right\|_F^2
\]
where $\hat{\mathbf{W}} \in \mathbb{R}^{3 \times d}$ stacks the $\ell_2$-normalised rows $\hat{\mathbf{w}}_V, \hat{\mathbf{w}}_A, \hat{\mathbf{w}}_D$.
This encourages the learned projection directions to be orthogonal, which is a necessary condition for applying steering vectors for individual PAD dimensions without cross-dimensional interference.

The full Phase~2 training objective is:
\[
    \mathcal{L} = \mathcal{L}_{\text{LM}}
    + \lambda_{\text{reg}}\,\mathcal{L}_{\text{reg}}
    + \lambda_{\text{con}}\,\mathcal{L}_{\text{con}}
    + \lambda_{\text{orth}}\,\mathcal{L}_{\text{orth}}
\]
with $\lambda_{\text{reg}} = 1.0$, $\lambda_{\text{con}} = 0.1$, $\tau = 1.0$, and $\lambda_{\text{orth}} = 0.1$.

\section{Training Strategy: Staged Unfreezing}

A further risk when warm-starting from an existing checkpoint is that the randomly initialised new heads produce large initial gradients that disturb the pre-trained latent representations before the heads have converged.
To mitigate this, the base language model is frozen for the first $N_{\text{freeze}} = 3$ epochs, during which only the three projection heads are updated.
After epoch~3, all parameters are unfrozen and trained jointly.
This staged unfreezing allows the heads to align with the existing latent geometry before the backbone is permitted to reorganise in response to the affective supervision signal.

\section{Phase 2 Results}

The Phase~2 architecture (CoconutGPT-Factored-v3) was trained for 35 epochs and evaluated on the same 861-sample EmoBank test split.
Table~\ref{tab:factored_v3} reports end-to-end generation metrics, and Table~\ref{tab:probe_v3} reports the linear probe results on the resulting latent representations.

\begin{table}[H]
\centering
\begin{tabular}{lcccc}
\toprule
\textbf{Model} & \textbf{Pearson V} & \textbf{Pearson A} & \textbf{Pearson D} & \textbf{Tol-Acc} \\
\midrule
CoconutGPT-Factored (v1, shared head) & 0.8466 & 0.8119 & 0.8411 & 45.30\% \\
\textbf{CoconutGPT-Factored-v3 (separate heads + orth.)} & 0.6857 & 0.5292 & 0.5890 & 43.44\% \\
\bottomrule
\end{tabular}
\caption{End-to-end generation metrics for the Phase~2 architecture. Despite the additional structural priors, separate heads with orthogonality regularisation \emph{degrade} generation performance on all three dimensions compared to the Phase~1 shared head.}
\label{tab:factored_v3}
\end{table}

\begin{table}[H]
\centering
\begin{tabular}{lccc}
\toprule
\textbf{Dimension} & \textbf{Pearson $r$ (test)} & \textbf{$R^2$ (test)} & \textbf{$R^2$ (train)} \\
\midrule
Valence (V)   & 0.69 & 0.42 & 0.93 \\
Arousal (A)   & 0.34 & $-$0.02 & 0.92 \\
Dominance (D) & 0.36 & 0.05 & 0.91 \\
\bottomrule
\end{tabular}
\caption{Linear probe on Phase~2 latent representations. Despite explicit per-dimension supervision, contrastive separation, and orthogonality regularisation, the linear decodability of A and D remains essentially unchanged from the Phase~1 baseline ($r \approx 0.37$).}
\label{tab:probe_v3}
\end{table}

Two observations stand out.
First, the linear probe results for A and D are statistically indistinguishable from the Phase~1 model: the orthogonality penalty did not produce a latent space in which Arousal and Dominance occupy separable linear directions.
Second, end-to-end generation accuracy declined --- Pearson A dropped from 0.81 to 0.53, and D from 0.84 to 0.59 --- suggesting that the additional structural constraints conflict with the model's natural encoding strategy rather than guiding it.
The freeze-then-unfreeze schedule and the additional regularisers appear to have over-constrained a 117M-parameter backbone that was already operating near its capacity for affective representation.

\section{Phase 3: Direct Steering Experiments}

To test whether the latent space of the Phase~1 best model (v2) admits causal manipulation through CAA-style activation addition~\cite{caa}, we computed per-dimension steering vectors and applied them at inference.

\subsection{Method}

For each PAD dimension $k \in \{V, A, D\}$, the 7{,}801 training samples were ranked by gold $k$ and split into a high-$k$ and low-$k$ contrast set ($k_{\text{frac}} = 0.2$, i.e.\ top/bottom 20\%).
Mean-pooled latent representations $\mathbf{h}_i = \frac{1}{|L_i|}\sum_{j \in L_i} \mathbf{e}_{i,j}$ were extracted at the latent token positions $L_i$ from \texttt{inputs\_embeds} after the model's multi-pass forward.
The unit-normalised steering vector for dimension $k$ is
\[
    \mathbf{d}_k = \frac{\bar{\mathbf{h}}^{\text{high}}_k - \bar{\mathbf{h}}^{\text{low}}_k}{\|\bar{\mathbf{h}}^{\text{high}}_k - \bar{\mathbf{h}}^{\text{low}}_k\|_2}.
\]
At inference, for a steering coefficient $\alpha \in \{-4, -2, -1, 0, +1, +2, +4\}$, we add $\alpha\,\mathbf{d}_k$ to every latent-token position of \texttt{inputs\_embeds} before greedy autoregressive decoding.
Predicted PAD values are parsed from the generated string.

\subsection{Results}

Table~\ref{tab:steering_summary} summarises the change in mean predicted PAD value as $\alpha$ ranges over $[-4, +4]$, alongside the cross-dimensional leakage onto the two non-target dimensions.

\begin{table}[H]
\centering
\begin{tabular}{lccc}
\toprule
\textbf{Steered dim} & \textbf{Target shift $\Delta\bar{k}$} & \textbf{Cross-V leak} & \textbf{Other leaks} \\
\midrule
V & $+0.097$ & --- & A: $+0.031$, D: $+0.022$ \\
A & $+0.112$ & $+0.012$ & D: $+0.006$ \\
D & $+0.027$ & $+0.037$ & A: $+0.027$ \\
\bottomrule
\end{tabular}
\caption{Mean predicted-PAD shift (1--5 scale) when sweeping $\alpha \in [-4, +4]$ along each unit steering vector. Direction is correct (monotonic) for all three dimensions, but the absolute magnitude is below the human-perceptible threshold, and steering D produces \emph{larger} interference on V than effect on D itself.}
\label{tab:steering_summary}
\end{table}

Three findings emerge from the steering sweep:
\begin{enumerate}
    \item \textbf{Direction is correct.} For all three dimensions, mean predicted PAD increases monotonically with $\alpha$, confirming that the contrastive mean-difference vector points in the intended direction.
    \item \textbf{Effect magnitude is below human-perceptible thresholds.} Even at $\alpha = +4$, the shift is at most $0.11$ on a 1--5 scale ($\approx$ 2.8\% of range). For comparison, the typical inter-annotator standard deviation on EmoBank is $\sigma \approx 0.4$, so the steering signal is roughly an order of magnitude smaller than annotation noise.
    \item \textbf{Cross-dimensional interference dominates Dominance.} Steering along $\mathbf{d}_D$ shifts $\bar{V}$ by $+0.037$ but only shifts $\bar{D}$ by $+0.027$. In other words, attempting to push the output toward higher Dominance changes Valence \emph{more} than Dominance itself.
    \item \textbf{Pearson $r$ is invariant to $\alpha$.} The per-sample ranking of predictions is unchanged across the entire $\alpha$ sweep; the steering operation acts as a near-uniform additive shift of the predicted distribution rather than as a per-sample modulation of affective content.
\end{enumerate}

\section{Discussion: Why VAD-Axis Steering Is Not Viable}

Taken together, the Phase~2 and Phase~3 results converge on a single conclusion: \textbf{the PAD axes do not constitute geometrically privileged, causally actionable directions in the latent space of a GPT-2-scale Coconut model}, even when the training objective explicitly rewards such alignment.
The implications for the original design of the Emotion Engine, in particular the proposed DJ-turntable interface in which users adjust V/A/D sliders, are significant:

\begin{itemize}
    \item Linear probe accuracy for A and D plateaus at $r \approx 0.37$ regardless of architecture (single shared head, separate heads, orthogonality regularisation, dimension-weighted contrastive loss). The information is present and \emph{decodable through the trained heads}, but it does not reside on a low-dimensional linear subspace that can be addressed by a CAA-style steering vector.
    \item Even when a steering direction is computed via mean-difference contrast --- the procedure that succeeds for sentiment and refusal in larger models~\cite{caa, hollinsworth-etal-2024-language} --- the resulting interventions are too weak to be perceptually meaningful and exhibit cross-dimensional interference that is comparable in magnitude to the target effect.
    \item A user moving a Dominance slider would observe predominantly a Valence change, violating the basic design promise of the interface (independent control of orthogonal affective axes).
\end{itemize}

We attribute this to two structural causes.
First, GPT-2's 768-dimensional residual stream and 12 layers may be too small a substrate to support three independent linear directions for affective control while simultaneously preserving language modelling capacity.
The success of CAA in the literature is consistently reported on much larger models (Llama-2-7B and above), where the residual stream has both the capacity and the over-completeness required for stable single-direction interventions~\cite{caa, reichman2025emotionsartthouunderstanding}.
Second, A and D are intrinsically correlated with V in natural text: high-arousal sentences are disproportionately negatively-valenced (anger, fear), and high-dominance sentences are disproportionately positively-valenced (pride, confidence).
Forcing the model to encode A and D along axes orthogonal to V conflicts with the empirical co-occurrence structure of the training data, which the model has every reason to internalise.

\section{Reorientation: From PAD Sliders to Data-Driven Affective Atoms}
\label{sec:reorientation}

Rather than continue trying to force the latent space into a fixed three-dimensional PAD frame, the next phase of the project will invert the problem: \emph{let the latent space dictate its own basis, and design the user interface around that basis}.
Concretely, we propose to replace the V/A/D sliders with a set of \textbf{data-driven affective atoms} discovered directly from the residual stream.

\subsection{Method Sketch}

\begin{enumerate}
    \item \textbf{Activation collection.} For each training sentence, collect the residual-stream activations $\mathbf{r}^{(\ell)} \in \mathbb{R}^{768}$ at a chosen layer $\ell$ (mid-network layers --- typically 6--9 of GPT-2 --- have been reported to carry the most semantically structured representations).
    \item \textbf{Decomposition.} Apply unsupervised decomposition --- PCA for orthogonal directions, NMF for non-negative, parts-based directions, or a sparse autoencoder for over-complete dictionaries --- to obtain a basis $\{\mathbf{b}_1, \dots, \mathbf{b}_K\}$ of affective atoms ($K \approx 8$--$32$).
    \item \textbf{Lexical anchoring.} For each atom $\mathbf{b}_j$, project the unembedding matrix $\mathbf{W}_U \in \mathbb{R}^{V \times 768}$ along $\mathbf{b}_j$ and report the top-$k$ tokens by inner product. The resulting human-readable label (e.g.\ ``\texttt{tender / hush / lullaby}'', ``\texttt{rage / strike / shatter}'') gives each atom an interpretable name without requiring it to align with a pre-specified psychological taxonomy.
    \item \textbf{Steering.} The user manipulates a small number of named atom intensities through sliders or knobs. The corresponding linear combination $\sum_j \alpha_j \mathbf{b}_j$ is added to the residual stream during generation, exactly as in CAA.
\end{enumerate}

\subsection{Why This Should Work Better}

The Phase~3 steering experiment failed because the target direction was \emph{externally specified} (the PAD axes from psychology) and only weakly aligned with directions of natural variation in the model's own representation.
Data-driven atoms, by construction, lie along directions of large variance in the residual stream and therefore correspond to interventions the model can actually express.
The trade-off is interpretability: rather than three psychologically pre-named axes, the user is presented with $K$ axes whose meaning must be inferred from their nearest-word projections.
This is a more honest interface: it surfaces the affective vocabulary the model has actually learned, rather than projecting it onto a frame that the model does not natively use.

\subsection{Revised Project Trajectory}

The final Emotion Engine will therefore consist of:
\begin{itemize}
    \item A frontend that displays $K$ atom sliders, each labelled by its top-3 nearest tokens and a 2D PCA scatter plot of training sentences for visual context.
    \item A backend that, given an input text and a vector of slider values, performs activation-addition steering at a fixed mid-layer of CoconutGPT-Factored-v2 (the best generation checkpoint from Phase~1) and returns the steered output.
    \item An evaluation protocol identical to the one outlined in Chapter~5, with the modification that ``Shift Accuracy'' is now measured per-atom rather than per-PAD-dimension.
\end{itemize}

This pivot preserves the project's core ambition --- pre-verbal, continuous, interactive emotional editing of generated text --- while abandoning the assumption, falsified by Phases~2 and~3, that the PAD frame is the right basis for that editing.

\chapter{Ethical Issues}

This project investigates an emotional AI engine that operates on the basis of latent representations of emotions. It will utilise recent concept-based reasoning paradigms (e.g.\ latent and implicit inference models) and use human judgement to evaluate the output. While the project does not involve invasive experiments or high-risk applications, it raises several ethical and legal considerations.

\section{Ethical Considerations}
The main ethical considerations relate to the involvement of human participants in the evaluation phase. Human evaluators are expected to assess the outputs of the model rather than disclose personal experiences or emotional states. Participation should be voluntary, and they should be free to withdraw at any time without consequence. The data collected shall be kept to a minimum (e.g.\ numerical evaluations and brief qualitative judgments only) and no personally identifiable information is requested.

The second ethical consideration concerns the nature of the system under study. This project explores the expression and regulation of emotions in AI systems. There is a potential ethical risk that these technologies, if misused, could be perceived as enabling emotional manipulation and deceptive emotional signals. However, the manipulation techniques addressed in this study are essentially identical to the widely adopted usage patterns of LLMs today. For example, instructions such as ``Please write more carefully'' or ``Could you rephrase this in a gentler manner?'' are already routinely provided as prompts, and the models adjust the emotional tone of their output accordingly. In other words, the significance of this project lies not in endowing models with new capabilities for emotional manipulation, but in visualising the structures of emotional expression already present within the models. The target of manipulation is not human emotion but the latent representations within the model. The value of this research lies in treating the ambiguous tone adjustments already occurring as a subject of direct study, rather than leaving them as a black box.

\section{Legal Considerations}

The project will also use open source software, pre-trained models (e.g.\ GPT-2) and public datasets. These resources will be used strictly in accordance with their respective licences and for non-commercial academic research purposes only. No proprietary, restricted or export-controlled software or information will be generated.

%% bibliography
\bibliography{main}
\bibliographystyle{unsrt}

\end{document}